\chapter{Physical AI Safety and the Hidden Observability Contract}
\label{ch:physical-ai-safety}
ORIUS is motivated by a simple but under-specified fact: modern physical AI systems
do not act on truth, they act on telemetry.  The deeper the software stack becomes,
the easier it is to hide the sensing contract that makes every downstream prediction,
optimization, and control action appear legitimate even when the observation channel is
degraded.  The manuscript therefore opens with a universal statement rather than a
witness-row statement: whenever action legality is evaluated on an observed state that
can diverge from the physical state, a hidden safety failure mode exists.

\section{Why physical AI safety is structurally different}
A physical AI stack does not end at prediction quality or planner confidence.  It ends
in a command that changes a battery dispatch, vehicle acceleration, plant input, alarm
state, guidance trajectory, or flight-control surface.  The danger is not merely that
the model can be wrong.  The danger is that the system can look internally coherent
while being externally unsafe because the software stack is reasoning over a degraded
observation surface.  This is what makes physical AI different from pure information
systems: mistakes are mediated by dynamics, latency, and irreversible physical
consequences rather than only by ranking error or classification loss.

This distinction also changes what counts as a defensible contribution.  A physical-AI
safety framework cannot rely only on better forecasting, better perception, or better
nominal control.  It must specify what happens when those upstream components are fed
stale, delayed, missing, or corrupted observations and then continue to act anyway.
ORIUS therefore begins not with task performance but with the hidden interface between
observation and action legality.

\section{Telemetry is the hidden control contract}
Every nominal controller implicitly signs a sensing contract.  It assumes that the state
estimate, telemetry packet, or inferred context is recent enough, synchronized enough,
and truthful enough for downstream legality checks to mean what the controller thinks
they mean.  Most deployed systems bury that contract inside preprocessing, filtering,
transport middleware, or state-estimation code.  Once it is buried, it becomes easy to
evaluate safety on the observed state and forget that the physical state may already have
drifted away.

That hidden contract is where the relevant literatures meet.  The uncertainty literature
contributes distribution-free predictive wrappers and coverage-aware calibration
surfaces \cite{vovk2005algorithmic,lei2018distribution,romano2019conformalized,gibbs2021adaptive,angelopoulos2023gentle,barber2023beyond,zaffran2022adaptive,stankeviciute2021conformal,xu2021dynamic,sesia2021comparison,bates2021distribution}.  Runtime assurance contributes supervisory architectures,
shield synthesis, and safety monitors that can veto or repair actions after a nominal
controller proposes them \cite{sha2001using,sha1998simplex,desai2019soter,bloem2015shield,bartocci2018introduction,havelund2004runtime,leucker2009brief,schierman2015run,bak2011runtime}.  Safe control contributes barrier methods,
robust model predictive control, and predictive safety filters that turn safety into a
computable constraint set rather than an ex post alarm \cite{ames2019control,ames2014cbf,wabersich2021predictive,nguyen2016exponential,wang2017safety,fisac2019general,mayne2005robust,langson2004tubes,rawlings2017mpc,ben2009robust,camacho2013model}.  Cyber-physical
systems research contributes the language of degraded telemetry, attack surfaces, change
detection, and monitored execution \cite{rajkumar2010cps,derler2012modeling,baheti2011cyberphysical,pasqualetti2013attack,teixeira2015secure,basseville1993detection,page1954continuous,page1955test,hinkley1971inference,liu2008isolation,breunig2000lof}.  ORIUS sits at the intersection of all
four because it treats telemetry degradation as a runtime action-semantics problem, not
merely as a detection or estimation nuisance.

\section{Why a universal safety layer is the right object}
The book uses the phrase \emph{physical AI} deliberately.  It refers to any stack in
which machine-learned or algorithmic decision logic is embedded in a loop that ends in
physical consequences: batteries, vehicles, industrial plants, ICU monitoring, guidance
systems, and aerospace control are all instances.  The universal claim of ORIUS is not
that one controller can solve them all.  The claim is that one defended safety layer can
express a common degraded-observation contract across them.

That is why ORIUS is written as a layer instead of as a replacement controller.
Domain-specific nominal intelligence remains domain-specific by design.  What can be
shared is the logic that decides when degraded observation widens the physically
plausible state set, tightens the admissible action set, forces a repair, triggers a
fallback, and records a certificate.  The reusable object is therefore not a universal
planner but a universal safety grammar.

\section{What the flagship ORIUS argument must establish}
That distinction matters for societal contribution.  The most credible path to broad
impact in physical AI is not another domain-specific optimizer, but a reusable,
auditable, and explicitly bounded safety layer that can sit between nominal intelligence
and real actuation.  ORIUS is written as that layer: a runtime contract that can be
instantiated repeatedly, promoted cautiously, and reviewed by people who care as much
about failure semantics as about nominal performance.

The opening chapters therefore do four jobs for the rest of the manuscript.  First, they name
the universal hazard: actions can be legal on observation while illegal in reality.
Second, they define the runtime object that addresses that hazard: a Detect--Calibrate--
Constrain--Shield--Certify kernel.  Third, they establish the claim boundary: ORIUS is
universal in architecture and runtime semantics, not automatically equal in empirical
maturity across all domains.  Fourth, they justify the monograph's universal-first
structure: one argument, one kernel, six domain chapters, and one visible parity gate
that prevents rhetoric from outrunning evidence.

The next chapter turns that intuition into a sharper organizing object.  It defines the
observation-action safety gap itself and explains why the entire book is strongest when
architecture-level universality and evidence-level promotion are kept distinct.
